\IfFileExists{stacks-project.cls}{%
\documentclass{stacks-project}
}{%
\documentclass{amsart}
}

% For dealing with references we use the comment environment
\usepackage{verbatim}
\newenvironment{reference}{\comment}{\endcomment}
%\newenvironment{reference}{}{}
\newenvironment{slogan}{\comment}{\endcomment}
\newenvironment{history}{\comment}{\endcomment}

% For commutative diagrams we use Xy-pic
\usepackage[all]{xy}

% We use 2cell for 2-commutative diagrams.
\xyoption{2cell}
\UseAllTwocells

% We use multicol for the list of chapters between chapters
\usepackage{multicol}

% This is generally recommended for better output
\usepackage{lmodern}
\usepackage[T1]{fontenc}

% For cross-file-references
\usepackage{xr-hyper}

% Package for hypertext links:
\usepackage{hyperref}

% For any local file, say "hello.tex" you want to link to please
% use \externaldocument[hello-]{hello}
\externaldocument[introduction-]{introduction}
\externaldocument[conventions-]{conventions}
\externaldocument[sets-]{sets}
\externaldocument[categories-]{categories}
\externaldocument[topology-]{topology}
\externaldocument[sheaves-]{sheaves}
\externaldocument[sites-]{sites}
\externaldocument[stacks-]{stacks}
\externaldocument[fields-]{fields}
\externaldocument[algebra-]{algebra}
\externaldocument[brauer-]{brauer}
\externaldocument[homology-]{homology}
\externaldocument[derived-]{derived}
\externaldocument[simplicial-]{simplicial}
\externaldocument[more-algebra-]{more-algebra}
\externaldocument[smoothing-]{smoothing}
\externaldocument[modules-]{modules}
\externaldocument[sites-modules-]{sites-modules}
\externaldocument[injectives-]{injectives}
\externaldocument[cohomology-]{cohomology}
\externaldocument[sites-cohomology-]{sites-cohomology}
\externaldocument[dga-]{dga}
\externaldocument[dpa-]{dpa}
\externaldocument[sdga-]{sdga}
\externaldocument[hypercovering-]{hypercovering}
\externaldocument[schemes-]{schemes}
\externaldocument[constructions-]{constructions}
\externaldocument[properties-]{properties}
\externaldocument[morphisms-]{morphisms}
\externaldocument[coherent-]{coherent}
\externaldocument[divisors-]{divisors}
\externaldocument[limits-]{limits}
\externaldocument[varieties-]{varieties}
\externaldocument[topologies-]{topologies}
\externaldocument[descent-]{descent}
\externaldocument[perfect-]{perfect}
\externaldocument[more-morphisms-]{more-morphisms}
\externaldocument[flat-]{flat}
\externaldocument[groupoids-]{groupoids}
\externaldocument[more-groupoids-]{more-groupoids}
\externaldocument[etale-]{etale}
\externaldocument[chow-]{chow}
\externaldocument[intersection-]{intersection}
\externaldocument[pic-]{pic}
\externaldocument[weil-]{weil}
\externaldocument[adequate-]{adequate}
\externaldocument[dualizing-]{dualizing}
\externaldocument[duality-]{duality}
\externaldocument[discriminant-]{discriminant}
\externaldocument[derham-]{derham}
\externaldocument[local-cohomology-]{local-cohomology}
\externaldocument[algebraization-]{algebraization}
\externaldocument[curves-]{curves}
\externaldocument[resolve-]{resolve}
\externaldocument[models-]{models}
\externaldocument[functors-]{functors}
\externaldocument[equiv-]{equiv}
\externaldocument[pione-]{pione}
\externaldocument[etale-cohomology-]{etale-cohomology}
\externaldocument[proetale-]{proetale}
\externaldocument[relative-cycles-]{relative-cycles}
\externaldocument[more-etale-]{more-etale}
\externaldocument[trace-]{trace}
\externaldocument[crystalline-]{crystalline}
\externaldocument[spaces-]{spaces}
\externaldocument[spaces-properties-]{spaces-properties}
\externaldocument[spaces-morphisms-]{spaces-morphisms}
\externaldocument[decent-spaces-]{decent-spaces}
\externaldocument[spaces-cohomology-]{spaces-cohomology}
\externaldocument[spaces-limits-]{spaces-limits}
\externaldocument[spaces-divisors-]{spaces-divisors}
\externaldocument[spaces-over-fields-]{spaces-over-fields}
\externaldocument[spaces-topologies-]{spaces-topologies}
\externaldocument[spaces-descent-]{spaces-descent}
\externaldocument[spaces-perfect-]{spaces-perfect}
\externaldocument[spaces-more-morphisms-]{spaces-more-morphisms}
\externaldocument[spaces-flat-]{spaces-flat}
\externaldocument[spaces-groupoids-]{spaces-groupoids}
\externaldocument[spaces-more-groupoids-]{spaces-more-groupoids}
\externaldocument[bootstrap-]{bootstrap}
\externaldocument[spaces-pushouts-]{spaces-pushouts}
\externaldocument[spaces-chow-]{spaces-chow}
\externaldocument[groupoids-quotients-]{groupoids-quotients}
\externaldocument[spaces-more-cohomology-]{spaces-more-cohomology}
\externaldocument[spaces-simplicial-]{spaces-simplicial}
\externaldocument[spaces-duality-]{spaces-duality}
\externaldocument[formal-spaces-]{formal-spaces}
\externaldocument[restricted-]{restricted}
\externaldocument[spaces-resolve-]{spaces-resolve}
\externaldocument[formal-defos-]{formal-defos}
\externaldocument[defos-]{defos}
\externaldocument[cotangent-]{cotangent}
\externaldocument[examples-defos-]{examples-defos}
\externaldocument[algebraic-]{algebraic}
\externaldocument[examples-stacks-]{examples-stacks}
\externaldocument[stacks-sheaves-]{stacks-sheaves}
\externaldocument[criteria-]{criteria}
\externaldocument[artin-]{artin}
\externaldocument[quot-]{quot}
\externaldocument[stacks-properties-]{stacks-properties}
\externaldocument[stacks-morphisms-]{stacks-morphisms}
\externaldocument[stacks-limits-]{stacks-limits}
\externaldocument[stacks-cohomology-]{stacks-cohomology}
\externaldocument[stacks-perfect-]{stacks-perfect}
\externaldocument[stacks-introduction-]{stacks-introduction}
\externaldocument[stacks-more-morphisms-]{stacks-more-morphisms}
\externaldocument[stacks-geometry-]{stacks-geometry}
\externaldocument[moduli-]{moduli}
\externaldocument[moduli-curves-]{moduli-curves}
\externaldocument[examples-]{examples}
\externaldocument[exercises-]{exercises}
\externaldocument[guide-]{guide}
\externaldocument[desirables-]{desirables}
\externaldocument[coding-]{coding}
\externaldocument[obsolete-]{obsolete}
\externaldocument[fdl-]{fdl}
\externaldocument[index-]{index}

% Theorem environments.
%
\theoremstyle{plain}
\newtheorem{theorem}[subsection]{Theorem}
\newtheorem{proposition}[subsection]{Proposition}
\newtheorem{lemma}[subsection]{Lemma}

\theoremstyle{definition}
\newtheorem{definition}[subsection]{Definition}
\newtheorem{example}[subsection]{Example}
\newtheorem{exercise}[subsection]{Exercise}
\newtheorem{situation}[subsection]{Situation}

\theoremstyle{remark}
\newtheorem{remark}[subsection]{Remark}
\newtheorem{remarks}[subsection]{Remarks}

\numberwithin{equation}{subsection}

% Macros
%
\def\lim{\mathop{\mathrm{lim}}\nolimits}
\def\colim{\mathop{\mathrm{colim}}\nolimits}
\def\Spec{\mathop{\mathrm{Spec}}}
\def\Hom{\mathop{\mathrm{Hom}}\nolimits}
\def\Ext{\mathop{\mathrm{Ext}}\nolimits}
\def\SheafHom{\mathop{\mathcal{H}\!\mathit{om}}\nolimits}
\def\SheafExt{\mathop{\mathcal{E}\!\mathit{xt}}\nolimits}
\def\Sch{\mathit{Sch}}
\def\Mor{\mathop{\mathrm{Mor}}\nolimits}
\def\Ob{\mathop{\mathrm{Ob}}\nolimits}
\def\Sh{\mathop{\mathit{Sh}}\nolimits}
\def\NL{\mathop{N\!L}\nolimits}
\def\CH{\mathop{\mathrm{CH}}\nolimits}
\def\proetale{{pro\text{-}\acute{e}tale}}
\def\etale{{\acute{e}tale}}
\def\QCoh{\mathit{QCoh}}
\def\Ker{\mathop{\mathrm{Ker}}}
\def\Im{\mathop{\mathrm{Im}}}
\def\Coker{\mathop{\mathrm{Coker}}}
\def\Coim{\mathop{\mathrm{Coim}}}

% Boxtimes
%
\DeclareMathSymbol{\boxtimes}{\mathbin}{AMSa}{"02}

%
% Macros for moduli stacks/spaces
%
\def\QCohstack{\mathcal{QC}\!\mathit{oh}}
\def\Cohstack{\mathcal{C}\!\mathit{oh}}
\def\Spacesstack{\mathcal{S}\!\mathit{paces}}
\def\Quotfunctor{\mathrm{Quot}}
\def\Hilbfunctor{\mathrm{Hilb}}
\def\Curvesstack{\mathcal{C}\!\mathit{urves}}
\def\Polarizedstack{\mathcal{P}\!\mathit{olarized}}
\def\Complexesstack{\mathcal{C}\!\mathit{omplexes}}
% \Pic is the operator that assigns to X its picard group, usage \Pic(X)
% \Picardstack_{X/B} denotes the Picard stack of X over B
% \Picardfunctor_{X/B} denotes the Picard functor of X over B
\def\Pic{\mathop{\mathrm{Pic}}\nolimits}
\def\Picardstack{\mathcal{P}\!\mathit{ic}}
\def\Picardfunctor{\mathrm{Pic}}
\def\Deformationcategory{\mathcal{D}\!\mathit{ef}}
\usepackage{cleveref}
\crefname{originalproposition}{original proposition}{original propositions}
\Crefname{originalproposition}{Original proposition}{Original propositions}
\makeatletter
\newlabel{proposition-quasi-compact-affine-diagonal}{{08DB}{source}{}{sourcepointer-1}{}}
\newlabel{proposition-quasi-compact-affine-diagonal@cref}{{[originalproposition][0][]08DB}{[1][0][]source}{}{}{}}
\makeatother
\crefname{originallemma}{original lemma}{original lemmas}
\Crefname{originallemma}{Original lemma}{Original lemmas}
\makeatletter
\newlabel{morphisms-lemma-affine-permanence}{{01SG}{source}{}{sourcepointer-2}{}}
\newlabel{morphisms-lemma-affine-permanence@cref}{{[originallemma][0][]01SG}{[1][0][]source}{}{}{}}
\makeatother
\crefname{originallemma}{original lemma}{original lemmas}
\Crefname{originallemma}{Original lemma}{Original lemmas}
\makeatletter
\newlabel{derived-lemma-adjoint-preserve-K-injectives}{{08BJ}{source}{}{sourcepointer-3}{}}
\newlabel{derived-lemma-adjoint-preserve-K-injectives@cref}{{[originallemma][0][]08BJ}{[1][0][]source}{}{}{}}
\makeatother
\crefname{originallemma}{original lemma}{original lemmas}
\Crefname{originallemma}{Original lemma}{Original lemmas}
\makeatletter
\newlabel{lemma-affine-pushforward}{{02MZ}{source}{}{sourcepointer-4}{}}
\newlabel{lemma-affine-pushforward@cref}{{[originallemma][0][]02MZ}{[1][0][]source}{}{}{}}
\makeatother
\title{The right derived functor of the pushforward agrees with the derived pushforward on the derived category of quasi-coherent modules}
\author{The Stacks Project Authors}
\date{stacks-project:perfect.tex:local-snapshot-20261012; GNU FDL 1.2 or later}
\begin{document}
\maketitle
\setcounter{subsection}{1}
\noindent\textit{Editorial note.} All delivered mathematical text is the author's own, copied from the pinned Stacks Project chapter file perfect.tex; no statement or proof step was rewritten, reordered or removed. The counted span is the author environment and its own proof environment. References inside the span are handed to the assembler as navigation stubs that carry the canonical tag of the cited tag only; whether each of them is load-bearing is recorded on the stub itself.\par\medskip
\noindent\textit{Allowed standard prerequisites (editorial).}\begin{itemize}
\item For a quasi-compact scheme with affine diagonal the derived category of quasi-coherent modules agrees with the full derived category of quasi-coherent sheaves. invoked at proof lines 1842-1846. Lets the proof replace one category by the other and compare the two functors on the same objects.
\item A pushforward along an affine morphism whose source category has an exact left adjoint preserves K-injective complexes, so it can be used to compute derived pushforwards. invoked at proof lines 1849-1862. Shows that pushing forward a K-injective complex of quasi-coherent modules from an affine open is again K-injective and represents its derived pushforward.
\item The pushforward along an affine morphism of a quasi-coherent complex represents its derived pushforward. invoked at proof lines 1858-1864. Identifies the values of both functors on the complexes coming from an affine open.
\end{itemize}\medskip
\ifdefined\hypertarget\hypertarget{sourcepointer-1}{}\fi
\noindent\textit{Source reference.} Original proposition 08DB: sources/stacks-perfect/pinned-original-source.tex, source lines 1801--1801. Cited at proof lines 1842-1845 for the equivalence of the categories; the comparison of the functors once they are identified is inside the counted span.\par\medskip
\ifdefined\hypertarget\hypertarget{sourcepointer-2}{}\fi
\noindent\textit{Source reference.} Original lemma 01SG: sources/stacks-chow/tags.txt, source lines 1--1. Used at proof lines 1856-1858 to see that the inclusion of an affine open and its composite with the given morphism are affine; the K-injectivity argument built on that is inside the counted span.\par\medskip
\ifdefined\hypertarget\hypertarget{sourcepointer-3}{}\fi
\noindent\textit{Source reference.} Original lemma 08BJ: sources/stacks-chow/tags.txt, source lines 1--1. Invoked at proof lines 1859-1862 to conclude that the pushforward from the affine open preserves K-injective complexes; the computation that follows is inside the counted span.\par\medskip
\ifdefined\hypertarget\hypertarget{sourcepointer-4}{}\fi
\noindent\textit{Source reference.} Original lemma 02MZ: sources/stacks-perfect/pinned-original-source.tex, source lines 1580--1580. Cited at proof lines 1863-1866 to identify the two functors on complexes pushed forward from the affine open; the generation argument extending the conclusion to all objects is inside the counted span.\par\medskip
\begin{lemma}
\label{lemma-direct-image-coherator}
Let $f : X \to Y$ be a morphism of schemes.
Assume $X$ and $Y$ are quasi-compact and have affine diagonal.
Then, denoting
$$
\Phi : D(\QCoh(\mathcal{O}_X)) \to D(\QCoh(\mathcal{O}_Y))
$$
the right derived functor of
$f_* : \QCoh(\mathcal{O}_X) \to \QCoh(\mathcal{O}_Y)$
the diagram
$$
\xymatrix{
D(\QCoh(\mathcal{O}_X)) \ar[d]_\Phi \ar[r] &
D_\QCoh(\mathcal{O}_X) \ar[d]^{Rf_*} \\
D(\QCoh(\mathcal{O}_Y)) \ar[r] &
D_\QCoh(\mathcal{O}_Y)
}
$$
is commutative.
\end{lemma}

\begin{proof}
Observe that the horizontal arrows in the diagram are
equivalences of categories by
Proposition \ref{proposition-quasi-compact-affine-diagonal}.
Hence we can identify these categories (and similarly for
other quasi-compact schemes with affine diagonal).
The statement of the lemma is that the canonical map
$\Phi(K) \to Rf_*(K)$ is an isomorphism for all $K$ in
$D(\QCoh(\mathcal{O}_X))$. Note that if $K_1 \to K_2 \to K_3 \to K_1[1]$
is a distinguished triangle in $D(\QCoh(\mathcal{O}_X))$ and
the statement is true for two-out-of-three, then it is true
for the third.

\medskip\noindent
Let $U \subset X$ be an affine open. Since the diagonal of $X$ is affine,
the inclusion morphism $j : U \to X$
is affine (Morphisms, Lemma \ref{morphisms-lemma-affine-permanence}).
Similarly, the composition $g = f \circ j : U \to Y$ is affine.
Let $\mathcal{I}^\bullet$ be a K-injective complex in $\QCoh(\mathcal{O}_U)$.
Since $j_* : \QCoh(\mathcal{O}_U) \to \QCoh(\mathcal{O}_X)$
has an exact left adjoint
$j^* : \QCoh(\mathcal{O}_X) \to \QCoh(\mathcal{O}_U)$
we see that $j_*\mathcal{I}^\bullet$ is a K-injective complex
in $\QCoh(\mathcal{O}_X)$, see
Derived Categories, Lemma \ref{derived-lemma-adjoint-preserve-K-injectives}.
It follows that
$$
\Phi(j_*\mathcal{I}^\bullet) =
f_*j_*\mathcal{I}^\bullet =
g_*\mathcal{I}^\bullet
$$
By Lemma \ref{lemma-affine-pushforward} we see that
$j_*\mathcal{I}^\bullet$ represents $Rj_*\mathcal{I}^\bullet$ and
$g_*\mathcal{I}^\bullet$ represents $Rg_*\mathcal{I}^\bullet$.
On the other hand, we have $Rf_* \circ Rj_* = Rg_*$.
Hence $f_*j_*\mathcal{I}^\bullet$ represents $Rf_*(j_*\mathcal{I}^\bullet)$.
We conclude that the lemma is true for any complex
of the form $j_*\mathcal{G}^\bullet$ with $\mathcal{G}^\bullet$
a complex of quasi-coherent modules on $U$. (Note that if
$\mathcal{G}^\bullet \to \mathcal{I}^\bullet$ is a quasi-isomorphism,
then $j_*\mathcal{G}^\bullet \to j_*\mathcal{I}^\bullet$ is a
quasi-isomorphism as well since $j_*$ is an exact functor
on quasi-coherent modules.)

\medskip\noindent
Let $\mathcal{F}^\bullet$ be a complex of quasi-coherent
$\mathcal{O}_X$-modules. Let $T \subset X$ be a closed subset
such that the support of $\mathcal{F}^p$ is contained in $T$
for all $p$. We will use induction on the minimal number $n$
of affine opens $U_1, \ldots, U_n$ such that
$T \subset U_1 \cup \ldots \cup U_n$. The base case $n = 0$ is trivial.
If $n \geq 1$, then set $U = U_1$ and denote $j : U \to X$ the
open immersion as above. We consider the map of complexes
$c : \mathcal{F}^\bullet \to j_*j^*\mathcal{F}^\bullet$.
We obtain two short exact sequences of complexes:
$$
0 \to \Ker(c) \to \mathcal{F}^\bullet \to \Im(c) \to 0
$$
and
$$
0 \to \Im(c) \to j_*j^*\mathcal{F}^\bullet \to \Coker(c) \to 0
$$
The complexes $\Ker(c)$ and $\Coker(c)$ are supported
on $T \setminus U \subset U_2 \cup \ldots \cup U_n$ and the result
holds for them by induction. The result holds for
$j_*j^*\mathcal{F}^\bullet$ by the discussion in the preceding
paragraph. We conclude by looking at the distinguished triangles
associated to the short exact sequences and using the initial
remark of the proof.
\end{proof}

\end{document}
